
\subsection{Proofs for Section~\ref{sec:timing}}\label{app:timing-proofs}

% ---------------------------------------------------------------------------
\subsubsection{Acquisition timing}\label{app:acquisition-proofs}

% compressed 2026-11-24: type (e): two sentences merged.
Word tests of either shape are allowed; three facts carry over from the
single-word class.

\emph{(T1) The tests see only $\theta$ and $\alpha$.} By
Lemma~\ref{lem:compile} and~\eqref{eq:word-joint}, whether a parameter
vector passes a word test of either shape, a calibration or a
probability-ratio test depends only on $\theta$ and~$\alpha$.

\emph{(T2) The register conditions are separate.} Since $\theta_m\ge\delta>0$
and $Q_m(g,h)\ge(1-\kappa)/|X|>0$, the positive entries of $P_m$ are those
of~$L_m$. Stochasticity, stationarity of~$\nu$, reciprocity and the
block-mass conditions concern $L$ alone.

% compressed 2026-11-24: (e) sentences of (T3) and Carrying merged, every
% check kept.
\emph{(T3) The projected-table modulus.} Since
$P_m((g,r),(h,r'))=\bar P_m(g,h)\,L_m(g,h;r,r')$ and $\bar P_m(g,h)$ sums
these entries over the $|R|$ values~$r'$,
$\|P_C-P_D\|_\infty<\varepsilon'$ gives
$\|\bar P_C-\bar P_D\|_\infty<|R|\varepsilon'$. If $C$ and $D$ have the same
register tables, with entries in $[0,1]$, then
$\|P_C-P_D\|_\infty\le\|\bar P_C-\bar P_D\|_\infty$. As in the proof of
Lemma~\ref{lem:transfer}(a), the row sums of $\bar P_m$ give $\theta_m$
and its sums weighted by $\chi_i(g)\chi_i(h)$ give $\theta_m\alpha_{mi}$,
so the projected tables determine $\theta$ and~$\alpha$.

\emph{Carrying.} Let $A$ and $C$ be feasible with $p_A\preceq p_C$, and
let $C^A$ have profile $p_A$, label probabilities $\theta_C$, eigenvalues
$\alpha_C$ and register tables~$L_A$. This configuration is feasible: the
model's conditions on $\theta$ and $\alpha$ hold at~$C$, and those on $L$
at~$A$; by~(T2), its positive entries form $\Gamma_{p_A}$ and its
decisions are those of~$A$. Every test held at $p_A$ is held at $p_C$ and
passed by~$C$, hence by $C^A$ via~(T1). The held word tests of $p_A$ are
admissible, as $A$ is feasible. Moreover $\bar P_{C^A}=\bar P_C$, $A^A=A$,
and $A_B=B^A$.

% compressed 2026-11-24: (c) the climb, written out a second time here, is
% replaced by a citation of the climb in the proof of Theorem directed,
% with the three facts it uses named.
\emph{Climbing.} If feasible $C$ and $D$ have $p_C\preceq p_D$ and the
same $\theta$ and~$\alpha$, then $D$ is finely reachable from~$C$. Indeed,
the climb in the proof of Theorem~\ref{thm:directed}, from $A'$ to~$B$,
applies with $C$ and $D$ in their places along the saturated chain from
$p_C$ to~$p_D$. That climb keeps $\theta$ and $\alpha$ fixed, so every
configuration built passes its held tests by~(T1). By~(T2), the registers
give feasible configurations, and the register interpolations use only
the affine argument of Lemma~\ref{lem:transfer}(d).

% compressed 2026-11-24: (e) sentences merged in (a)-(c); both directions
% of (a) and (b) kept.
\begin{proof}[Proof of Theorem~\ref{thm:projected-components}]
(a) Carrying shows that $A_B$ is feasible. If $A_B$ is finely reachable
from~$A$, append a climb from $A_B$ to~$B$. Conversely, given rational
$\varepsilon>0$, put $\varepsilon'=\varepsilon/(|R|+1)$ and take a
reconfiguration sequence $A=C_0,C_1,\dots,C_n=B$ at
tolerance~$\varepsilon'$. Profiles never decrease along it, so every
$C_i^A$ is feasible at~$p_A$. A profile step keeps $\theta$ and $\alpha$,
so $C_{i-1}^A=C_i^A$; a numerical step gives carried tables closer than
$|R|\varepsilon'<\varepsilon$, by both bounds of~(T3). Deleting repetitions
leaves numerical steps at~$p_A$, at tolerance~$\varepsilon$, from $A$
to~$A_B$.

(b) Steps never lower the profile, so fine reachability requires
$p_A\preceq p_B$. Under this condition, by~(a) and
$\bar P_{A_B}=\bar P_B$, it suffices that $A_B$ is finely reachable from
$A$ exactly when $\bar P_A$ and $\bar P_B$ share a component of
$\mathcal K_p$, where $p=p_A$. The set $H_p$
of projected tables of feasible rational configurations at~$p$ is dense
in~$\mathcal K_p$, which is closed and bounded, hence compact. Every
sequence from $A$ to $A_B$ stays at~$p$ (Remark~\ref{rem:inputs}(ii)). %
If $A_B$ is finely reachable from~$A$, then for each $\varepsilon>0$ take a
sequence at a positive rational tolerance $\varepsilon'$ with
$|R|\varepsilon'<\varepsilon$. By~(T3), the projected sequence lies in
$H_p$ and has steps below~$\varepsilon$. Lemma~\ref{lem:compact-chain}
therefore puts $\bar P_A$ and $\bar P_B$ in one component. Conversely, given rational $\varepsilon>0$,
Lemma~\ref{lem:compact-chain} gives $\bar P_A=h_0,h_1,\dots,h_k=\bar P_B$
in $H_p$ with consecutive sup distances below~$\varepsilon$. Each $h_j$ is
the projected table of a feasible rational $C_j$ at~$p$. The carried
$C_j^A$ keep these tables and share one register, so by~(T3) consecutive
full tables differ by less than~$\varepsilon$. As the
projected tables determine $\theta$ and~$\alpha$, the first is $A$ and the
last~$A_B$.

(c) Let $C$ be feasible and rational at~$q$, and $A$ feasible and rational
at~$p$, which is attained. Carrying gives $C^A$ at~$p$ with
$\bar P_{C^A}=\bar P_C$, so $H_q\subseteq H_p$ and the closures are
nested.

(d) Apply~(b) to $A$ and $B$, and to $A^+$ and~$B$. The order conditions
hold because $p_A\preceq p_{A^+}\preceq p_B$.
\end{proof}

% compressed 2026-11-24: (d) the reading of (d) is cut to the one fact it
% needs; the sign-example justification is kept.
A feasible acquisition keeps the parameter vector, so
$\bar P_{A^+}=\bar P_A$ in~(d). In the sign example, the feasible
configurations at $p_A$ have label probabilities $1/2$, $z=1/4$ and any
rational $y$ with $|y|\le\kappa$. Their projected tables are affine
in~$y$, so $\mathcal K_{p_A}$ is a segment, whereas at $p_{A^+}$ the test
$\mathsf{yy}\sim\mathsf{zz}$ leaves only $y=\pm1/4$.

% compressed 2026-11-24: (e) sentences merged in both proofs; every check
% kept.
\begin{proof}[Proof of Theorem~\ref{thm:acquisition}]
By Theorem~\ref{thm:directed}, $B$ is finely reachable from $A$ exactly
when $p_A\preceq p_B$ and $\SC_{p_A}(\alpha_A)=\SC_{p_A}(\alpha_B)$, and
from $A^+$ exactly when $p_{A^+}\preceq p_B$ and
$\SC_{p_{A^+}}(\alpha_{A^+})=\SC_{p_{A^+}}(\alpha_B)$. Both order
conditions hold, and $\alpha_{A^+}=\alpha_A$. The profile comparisons, the
recovery of parameters from supplied tables and the compilation of
$\Phi_{p_A}$ and $\Phi_{p_{A^+}}$ take polynomial time
(Appendix~\ref{app:tractable-proofs}). The test makes two signed-core
comparisons: using Theorem~\ref{thm:certificates}(a), it computes the
cores of $\supp\alpha_A$ and $\supp\alpha_B$ under $\Phi_{p_A}$ and then
under $\Phi_{p_{A^+}}$, and compares signs on them.
\end{proof}

\begin{proof}[Proof of Proposition~\ref{prop:source-schedule}]
Fix rational $\varepsilon>0$. The proof of
Theorem~\ref{thm:projected-components}(a) turns a sequence from $A$ to $B$
at tolerance $\varepsilon/(|R|+1)$ into one from $A$ to $A_B$ at
tolerance~$\varepsilon$ at~$p_A$. For $1\le j\le k$, let $C_j$ be $B$
carried to~$q_j$ with the register of some feasible configuration
at~$q_j$; it is feasible as $q_j\preceq p_B$. With $C_0=A_B$ and
$C_{k+1}=B$, consecutive configurations have ordered profiles and the
parameters $\theta_B$ and $\alpha_B$, so climbs at tolerance $\varepsilon$
join them, changing only registers and profiles. If
$\Gamma_{q_j}=\Gamma_{p_B}$, use the register of $B$ at~$q_j$: its support
is right, each decision of $q_j$ is one of~$p_B$, as the partitions of
$q_j$ are coarser, and the tests held at $q_j$ are held at~$p_B$. Then
$C_j$ has the full parameter vector of~$B$.
\end{proof}

% ---------------------------------------------------------------------------
\subsubsection{Order properties of the core}\label{app:core-order-proofs}

\begin{proposition}[Empty cores]\label{prop:empty-core}
Fix a budgeted signed-binomial system on~$J$, and $S\subseteq J$.
(a)~$\core(S)=\emptyset$ if and only if $S$ has no nonempty irreducible
subset. (b)~If some weight vector with $\Delta_S\eta=0$ is positive on all
of~$S$, in particular if every test active on $S$ has equal total degrees
on its two sides, then one peeling step leads from $S$ to~$\emptyset$.
(c)~Any two feasible rational points with empty cores are finely
connected, with no boundedness assumed.
\end{proposition}

% compressed 2026-11-24: (e) sentences merged in (a)-(c).
\begin{proof}[Proof of Proposition~\ref{prop:empty-core}]
(a) By Lemma~\ref{lem:irreducible}, $\core(S)$ is an irreducible subset
of $S$ containing every irreducible subset of~$S$, so it is empty exactly
when no irreducible subset of $S$ is nonempty.
(b) A weight vector positive on all of $S$ defines the step
$S\to\emptyset$. If every test $x^a=c\,x^b$ active on $S$ has
$\sum_ja_j=\sum_jb_j$, then $\eta=\mathbf 1_S$ gives $(a-b)\cdot\eta=0$
for each, as active tests have exponents only in~$S$.
(c) If $x$ is feasible and rational with $\core(\supp x)=\emptyset$,
iterating Lemma~\ref{lem:restriction}(a), which uses no boundedness,
along a peeling of $\supp x$ to $\emptyset$ finely connects $x$ to the
feasible point~$0$. Do the same for the other point; fine connectivity is
an equivalence relation.
\end{proof}

% compressed 2026-11-24: (e) sentences merged; the degree count of each
% equation type, both directions of the core claim, the order condition
% and the feasibility of A_0 are kept.
\begin{proof}[Proof of Theorem~\ref{thm:no-calibration}]
Let $\alpha$ satisfy the equations of $\Phi_p$, and put $S=\supp\alpha$. A
reversal equation has degree one on each side, the probe equations
$\alpha_i(w)=\alpha_i(w')$ of a word test $w\sim w'$ have degrees
$|w|=|w'|$, and a calibration $\alpha_{mi}=c$ has degrees one and zero. If
$c=0$, then $\alpha_{mi}=0$, so $(m,i)\notin S$ and the calibration is
active on no subset of~$S$. Hence, if $p$ holds no nonzero calibration,
every test active on $S$ has equal degrees, and
Proposition~\ref{prop:empty-core}(b) peels $S$ to $\emptyset$ in one step,
so $\SC_p(\alpha)=0$. Conversely, if $p$ holds $\alpha_{mi}=c\neq0$, then
$(m,i)\in S$. On every support containing $(m,i)$, that calibration is
active; its row is the unit vector at $(m,i)$, forcing $\eta_{mi}=0$.
Hence no peeling step removes $(m,i)$, and $\SC_p(\alpha)_{mi}=\sgn c\neq0$.

Now let $A$ be feasible with no nonzero calibration held at~$p_A$, and $B$
feasible. Fine reachability needs $p_A\preceq p_B$, as steps never lower
the profile. If $p_A\preceq p_B$, then $p_B$ holds the tests of $p_A$, so
$\alpha_B$ satisfies the equations of $\Phi_{p_A}$
(Lemma~\ref{lem:separate}(b)); both signed cores at $p_A$ are zero, and
Theorem~\ref{thm:directed} gives fine reachability. At a fixed profile
with no nonzero calibration the order condition holds both ways, so any
two feasible configurations there are mutually finely reachable. Finally,
let $A_0$ have the profile, label probabilities and register tables of~$A$
and all eigenvalues zero. The zero vector satisfies the probe equations of
the held word tests (whose words are nonempty), the reversal equations,
the held calibrations (all zero), and the budgets. With the label
probabilities and register tables of~$A$, $A_0$ is thus feasible by
Lemma~\ref{lem:separate}, and mutually finely reachable with $A$ by the
fixed-profile claim just proved.
\end{proof}

% compressed 2026-11-24: (d) the first example's restatement of the body
% (eta=(1,1) for y^2=z^2) is cut to its feasibility check; (e) merges.
\emph{The examples after Theorem~\ref{thm:no-calibration}.} Without the
calibration, the sign example's source $y=z=1/4$ stays feasible. For
$(xy)^2=xy$ with $|x|+|y|\le2$, both $(1,1)$ and $(-1,-1)$ are feasible.
On $\{x,y\}$ the balance $2\eta_x+2\eta_y=\eta_x+\eta_y$ forces $\eta=0$,
so $\{x,y\}$ is irreducible; hence the signed cores $(+,+)$ and $(-,-)$
differ, and by Theorem~\ref{thm:core} the points are not finely connected.
With $z=0$ held and $\mathsf{yy}\sim\mathsf{zz}$
not held, the row of $z=0$ is active on $\{z\}$ and forces $\eta_z=0$, so
the formal support $\{z\}$ is irreducible, yet no eigenvalue vector
passing the calibration has that support. The all-zero eigenvalue vector
shows that the profile is attained.

\begin{corollary}[Adding tests refines the classes]\label{cor:refine}
(a)~Let every test of a budgeted signed-binomial system $\Phi$ occur in
$\Phi'$, with the same budgets. Then, for every~$S$,
$\core_\Phi(S)\subseteq\core_{\Phi'}(S)$ and
\[
  \core_\Phi(\core_{\Phi'}(S))=\core_\Phi(S)=\core_{\Phi'}(\core_\Phi(S)),
\]
equal $\Phi'$-signed cores imply equal $\Phi$-signed cores,
and points finely connected under $\Phi'$ are finely connected
under~$\Phi$.
(b)~In the single-word class, if $A$ and $A'$ are feasible with
$p_A\preceq p_{A'}$ and $\alpha_A=\alpha_{A'}$, every configuration finely
reachable from $A'$ is finely reachable from~$A$. In particular an
acquisition never enlarges the set of configurations finely reachable from
its source.
\end{corollary}

\begin{proof}[Proof of Corollary~\ref{cor:refine}]
(a) A test of $\Phi$ active on $S$ is a test of $\Phi'$ active on~$S$, so
every $\Phi'$-peeling step is a $\Phi$-peeling step with the same weights.
Hence $S$ peels under $\Phi$ to $\core_{\Phi'}(S)$, and
Lemma~\ref{lem:restriction}(c) gives
$\core_\Phi(S)\subseteq\core_{\Phi'}(S)$ and
$\core_\Phi(\core_{\Phi'}(S))=\core_\Phi(S)$.
Also \(\core_{\Phi'}(\core_\Phi(S))=\core_\Phi(S)\). Deflation gives one
inclusion, and idempotence and monotonicity in the test list give
\(\core_\Phi(S)=\core_\Phi(\core_\Phi(S))\subseteq\core_{\Phi'}(\core_\Phi(S))\).
Both absorption identities also hold for the signed-core maps on sign
vectors, since each composition keeps the original signs on the resulting
support. If
$\SC_{\Phi'}(x)=\SC_{\Phi'}(x')$, then $\supp x$ and $\supp x'$ have one
$\Phi'$-core $C'$, on which $x$ and $x'$ have the same signs, and both
$\Phi$-cores equal $\core_\Phi(C')\subseteq C'$, so
$\SC_\Phi(x)=\SC_\Phi(x')$. A sequence of feasible rational points
of~$\Phi'$ is one of~$\Phi$, with the same steps, as the budgets agree.
(b) If $B$ is finely reachable from $A'$, Theorem~\ref{thm:directed}
gives $p_{A'}\preceq p_B$, hence $p_A\preceq p_B$, and
$\SC_{p_{A'}}(\alpha_{A'})=\SC_{p_{A'}}(\alpha_B)$. The held tests are
nested, so every equation of $\Phi_{p_A}$ occurs in $\Phi_{p_{A'}}$, with
the same budgets. By~(a) and $\alpha_A=\alpha_{A'}$,
$\SC_{p_A}(\alpha_A)=\SC_{p_A}(\alpha_B)$, and Theorem~\ref{thm:directed}
gives fine reachability from~$A$.
\end{proof}

% compressed 2026-11-24: (d) the sign-example paragraph that followed this
% proof is cut; it restated the sign example of Section 4.5 (Phi_{p_A}
% gives only eta_z=0; after the acquisition {y,z} is irreducible).

\begin{proposition}[Interior and closure operators]\label{prop:interior}
Fix a budgeted signed-binomial system on~$J$. The core is an interior
operator on subsets of~$J$: monotone, with $\core(S)\subseteq S$ and
$\core(\core(S))=\core(S)$. The complement of the core on removed
coordinates, $Z\mapsto J\setminus\core(J\setminus Z)$, is a closure
operator whose closed sets are the complements of irreducible supports.
At a profile $p$ of a single-word model, if a feasible configuration at~$p$
has eigenvalue support~$S$, some feasible configuration at~$p$ has
eigenvalue support $\core_p(S)$, the greatest irreducible such support
inside~$S$.
\end{proposition}

% compressed 2026-11-24: (e) sentences merged; (d) the closing
% sign-example sentence is cut, as it restated Section 6.2 and the sign
% example of Section 4.5 (Phi_{p_A} gives only eta_z=0, so y peels).
\begin{proof}[Proof of Proposition~\ref{prop:interior}]
Peeling only removes coordinates, so $\core(S)\subseteq S$, and $S$ peels
to $\core(S)$, which therefore has core $\core(S)$, by
Lemma~\ref{lem:restriction}(c). If $U\subseteq S$, then $U$ peels to $U\cap\core(S)$
by Lemma~\ref{lem:restriction}(b), so $\core(U)\subseteq\core(S)$.
Complementation in $J$ reverses inclusion and is an involution, so
$Z\mapsto J\setminus\core(J\setminus Z)$ is monotone, idempotent and
extensive, and its fixed sets are the $Z$ with
$\core(J\setminus Z)=J\setminus Z$, that is, with $J\setminus Z$
irreducible.

For the last claim, let $A$ be feasible at~$p$ with $S=\supp\alpha_A$.
Iterating Lemma~\ref{lem:restriction}(a) along a peeling of $S$ to
$\core_p(S)$ gives a feasible rational point of $\Phi_p$ with support
$\core_p(S)$ that agrees with $\alpha_A$ there. With $\theta_A$ and $L_A$,
this point forms a feasible configuration at~$p$
(Lemma~\ref{lem:separate}). By Lemma~\ref{lem:irreducible}, $\core_p(S)$
contains every irreducible subset of~$S$, attained or not.
\end{proof}
